\documentclass[a4paper,12pt,oneside]{article}
\makeatletter
\def\input@path{{../../}}
\makeatother
\usepackage{discrete-practice}

% Источник условий: https://nebulous-name-0d3.notion.site/1-20-02-2024-23-59-95a99f44fd74446b856048489b454f8b?pvs=25
\begin{document}

\HomeworkHeader{2}{1}{Графы}

\Task{Задание 1}
Граф $G$ задан матрицей смежности:
\[
  \begin{pmatrix}
    0 & 1 & 1 & 0 & 0 \\
    1 & 0 & 1 & 0 & 0 \\
    1 & 1 & 0 & 1 & 1 \\
    0 & 0 & 1 & 0 & 1 \\
    0 & 0 & 1 & 1 & 0
  \end{pmatrix}.
\]
\begin{itemize}
  \item Нарисуйте этот граф.
  \item Постройте его матрицу инцидентности.
  \item Есть ли в нём висячие вершины?
  \item Является ли он связным?
  \item Ответьте на те же вопросы для дополнительного графа.
\end{itemize}

\Task{Задание 2}
Для графа $G$ и дополнительного к нему постройте:
\begin{alphlist}
  \item любой частичный граф;
  \item надграф;
  \item подграф;
  \item нуль-граф.
\end{alphlist}

\Task{Задание 3}
Нарисуйте граф для каждого списка степеней его вершин:
\begin{alphlist}
  \item $1,1,3,3,4,4,6,6,2$;
  \item $2,2,2,3,3,4,5,5$.
\end{alphlist}
Является ли он двудольным графом?

\Task{Задание 4}
Сколько рёбер в дополнительном графе к простому неориентированному графу $G$, если в $G$:
\begin{alphlist}
  \item степени вершин равны $1,2,2,3,3$?
  \item степени вершин равны $7,7,5,4,4,4,6,4,3$?
  \item число вершин равно $12$, а число рёбер равно $23$?
\end{alphlist}

\clearpage

\Task{Задание 5}
Полагаем, что граф простой.
\begin{alphlist}
  \item Какое наименьшее число рёбер может быть в связном графе с $n$ вершинами?
  \item Какое наибольшее число рёбер может быть в несвязном графе с $n$ вершинами?
  \item Какое наименьшее число рёбер может быть в графе с $n$ вершинами и $k$ компонентами связности?
\end{alphlist}

\Task{Задание 6}
Найдите наименьшее и наибольшее число компонент связности в простом неориентированном графе, если в нём $14$ вершин и $5$ рёбер.

\Task{Задание 7}
Найдите наименьшее и наибольшее количество вершин в простом неориентированном графе:
\begin{alphlist}
  \item если в нём $22$ ребра и $3$ компоненты связности;
  \item если в нём $132$ ребра и $12$ компонент связности.
\end{alphlist}

\Task{Задание 8}
Перечислите компоненты сильной связности. Постройте конденсацию графа.
\begin{center}
  \includegraphics[width=.43\linewidth]{HW_1_graph.png}
\end{center}

\end{document}
